\documentclass[11pt]{article}
\usepackage{explorer}
\newcommand{\ExplorerLanguage}{EN}
\newcommand{\ExplorerNumber}{02}
\newcommand{\ExplorerKind}{Worked solutions}
\begin{document}
\ExplorerTitle{High school · olympiad thinking}{The pigeonhole principle}
\ExplorerSection{Exercise 1}
Show that among 25 people, at least three were born in the same month.\par
At most two per month would give at most 24 people. Therefore one month contains at least three.\par
\ExplorerSection{Exercise 2}
Choose six distinct numbers from \{1,…,10\}. Prove that two sum to 11.\par
The five boxes are \{1,10\}, \{2,9\}, \{3,8\}, \{4,7\}, \{5,6\}. Six choices fill at least one pair.\par
\ExplorerSection{Exercise 3}
Prove that among eight integers, two have a difference divisible by 7.\par
There are seven possible remainders modulo 7. Two integers share a remainder, so their difference is a multiple of 7.\par
\ExplorerSection{Exercise 4}
Five points lie in a square of side 2, including its boundary. Show that two are at distance at most \(\sqrt{2}\).\par
Divide the square into four unit squares, assigning each boundary point to just one square. Two points share a unit square, whose diagonal is \(\sqrt{2}\).\par
\ExplorerSection{Problem: complete reasoning}
Choose 11 distinct integers from 1 to 20. Write each as 2ᵏm with m odd and \(k \geq  0\). How many values of m are possible? Prove that one selected integer divides another. Show that the conclusion can fail with only ten integers.\par
Remove all factors of 2. The ten possible odd parts are 1,3,…,19. Two chosen numbers share an odd part: 2ᵃm and 2ᵇm, with \(a < b\) because they are distinct. The second is 2ᵇ⁻ᵃ times the first. For ten numbers choose 11,12,…,20: twice any of them exceeds 20, so none divides another distinct one.\par
\end{document}
